\scanpage{91}
et en considérant la courbe $X'$, ayant $I'$ comme ensemble de pts à l'infini, et $A'$ comme ens. des pts doubles, obtenue en ``bouchant les trous'' relatifs aux $i\in I\setminus I'$, \deleted{\unsure{et|en}} et en normalisant $X$ en les pts de $A\setminus A'$. On trouve donc $X'=\widehat X'\setminus I'$, $\widehat X'$ étant la normalisée de $\widehat X$ relatif aux pts de $A\setminus A'$. On fera attention que $\overset{\deleted{\wedge}}{X}'$ \uline{ne s'envoie pas} dans $\overset{\deleted{\wedge}}{X}$ (sauf si $I'=I$), mais seulement dans $\widehat X$.

\deleted{On peut aussi aller plus loin, et} La construction de $X'$ en termes de $X$ est donc complète, dès qu'on a donné des parties $I'\subset I$ et $A'\subset A$ --- cela revient, en termes du \deleted{la maquette} graphe associé à $X$, à se donner

\noindent a) un sous-graphe \inserted{$G'$} de $G$ ayant $=$ ensemble des sommets

\noindent b) un sous-ens. $I'$ de $I$.

\deleted{Mais on} Le graphe de $X'$ est donc le sous-graphe en question. Le \deleted{MD}-maquette
\margin{90}
